\section*{Erd\H{o}s Problem 114: lengths of polynomial lemniscates}
\subsection*{Statement and scope}
For a monic degree-$n$ polynomial $p$, is the length of $\{z:|p(z)|=1\}$ maximal at $p(z)=z^n-1$? The assertion for every $n$ is unresolved in this write-up. We prove the maximization within the entire translated binomial family and compute its value, including the constant term in its asymptotic expansion.

This develops the valid length integral from the earlier GPT-5.2 attempt. Current statement: \url{https://www.erdosproblems.com/114}. Tao proves the full assertion for all sufficiently large degrees in \url{https://arxiv.org/abs/2512.12455}; that theorem is context, not an input to the restricted proof below. No novelty is claimed.

\subsection*{The inverse-image length formula}
Translation of the variable does not affect length, so consider $p(z)=z^n-c$. For $n>1$ put $s=1-1/n\in(0,1)$. On the level curve, $z^n=c+e^{it}$. Away from the finitely many singular parameters, each of the $n$ inverse branches has speed
\[
 \frac1n |c+e^{it}|^{1/n-1}.
\]
Summing branch lengths gives
\[
 L_n(c)=\int_0^{2\pi}|c+e^{it}|^{-s}\,dt. \tag{1}
\]
This counts each smooth arc once; the branch singularity when $|c|=1$ is integrable because $s<1$. Individual branches need not return to their initial values after one circuit; summing all branches is what justifies (1).

Rotation of $t$ shows that (1) depends only on $r=|c|$. For $0\le r<1$, the binomial series and termwise integration, justified by absolute uniform convergence for each fixed $r<1$, give
\[
 L_n(r)=2\pi\sum_{j=0}^\infty
 \left(\frac{(s/2)_j}{j!}\right)^2 r^{2j}. \tag{2}
\]
Indeed expand $(1+r e^{it})^{-s/2}$ and its conjugate; only equal indices survive integration. All coefficients are positive, so $L_n(r)$ strictly increases on $[0,1)$.

For $r>1$, factoring out $r$ gives
\[
 L_n(r)=r^{-s}L_n(1/r)
 =2\pi\sum_{j=0}^\infty
 \left(\frac{(s/2)_j}{j!}\right)^2 r^{-s-2j}, \tag{3}
\]
which strictly decreases. Continuity at $r=1$ follows by dominated convergence: for $r\in[1/2,3/2]$,
\[
 |r+e^{it}|^2=(r-1)^2+4r\cos^2(t/2)\ge2\cos^2(t/2),
\]
and $|\cos(t/2)|^{-s}$ is integrable. Equations (2)--(3) prove that the unique maximizing modulus is $|c|=1$. Thus the proposed extremizer is correct throughout the family $(z-a)^n-c$.

\subsection*{Exact length and a justified asymptotic}
At $c=1$, (1) yields
\[
 L_n(1)=2^{1/n}B\left(\frac1{2n},\frac12\right), \tag{4}
\]
where $B$ is the beta integral. For $\epsilon\downarrow0$,
\[
 B(\epsilon,\tfrac12)-\frac1\epsilon
 =\int_0^1 t^{\epsilon-1}\big((1-t)^{-1/2}-1\big)\,dt
 =2\log2+O(\epsilon). \tag{5}
\]
The limit integral is $2\log2$, by $u=\sqrt{1-t}$. The error estimate follows from $|t^\epsilon-1|\le\epsilon|\log t|$ and the integrability of
$t^{-1}|\log t|((1-t)^{-1/2}-1)$ at both endpoints. Combining (4)--(5) with
$2^{1/n}=1+(\log2)/n+O(n^{-2})$ gives
\[
 L_n(1)=2n+4\log2+O(n^{-1}).
\]
For $n=1$, every monic linear polynomial has a unit-circle lemniscate of length $2\pi$.

\subsection*{Remaining gap}
The positive-coefficient comparison uses the two-term form of the polynomial. It supplies no upper bound for a general monic polynomial in the degrees not covered by the cited large-degree theorem. That is the missing comparison for the full question.
\subsection*{COMPLETION ESTIMATE}
\textbf{COMPLETION: 45\%}
